\cotefile{29}
\scanpage{29}
% archive sheet number: 28

est``l'image inverse'' de $\underline{R}_{(S)}$, toujours sous les conditions de l'exemple 3. Alors il est \uline{trivial} que $\operatorname{Rev}^{\underline{R}_{(S)}}(S)\qquad\operatorname{Rev}^{\underline{R}_{(T)}}(T)$ est une équivalence, utilisant le résultat analogue pour les revêtements étales de schémas $X$ et $X\unsure{x|\times}_{S}T$ (où $X\ \deleted{\unsure{\mathrm{est}}}=\underline{s}(M)$ est un schéma \deleted{\unsure{cxxvx|connx}} sur $S$ ..). En particulier, quand on dispose d'un point géométrique non ramifié de $T$, on obtient un isomorphisme sur les groupes fondamentaux. Le même raisonnement, après localisation étale à peu près évidente, s'applique au\deleted{\unsure{x|s}} cas kummérien \deleted{\unsure{a)}}.

J'ai fait allusion p\deleted{\unsure{h}}lusieurs fois à la donnée de définir une donnée de ramification image inverse d'une donnée de ramification sur $S$ par un morphisme $T\quad S$; la définition explicite serait nécessairement un peu lourde; je me borne à signaler que dans le cas où $\underline{R}$ est défini par $(Z,G)$ comme dans l'exemple 1, son image inverse est définie par $(Z_T=Zx_ST,G)$. L'ennui avec cette définition générale, c'est qu'elle ne s'applique pas directement au cas où $\underline{R}$ est une donnée kummérienne associée à un systême $(D_i)$, car la donnée image inverse au sens précis précédent n'\uline{est pas} la donnée correspondante au\deleted{\unsure{xxRx|xxex}} systême $(D_{i\,T})$ même à équivalence de catégories fibrées près; elle lui est seulement ``essentiellement équivalente'' au sens du n°1. Seule une rédaction en forme permettra de se rendre compte s'il y a lieu d'introduire cette notion d'image inverse, ou s'il ne vaut pas mieux donner des conditions axiomatiques sur la donnée $(\unsure{*|\pi})$, sous les conditions de l'exemple 3, qui assurent moralement que $\underline{R}_{(T)}$ est essentiellement équivalente à l'image inverse de $\underline{R}_{(S)}$. Bien sûr, ces conditions, pour \deleted{\unsure{avoir|avxix}} être utiles, doivent s'appliquer directement aux cas kummériens. Ce sera sans doute le parti le plus raisonnable. Il me semble également qu' serait raisonnable, assez rapidement après les définitions général-
